\section{Finite acquisition through singular nonreset dynamics}\label{sec:singular}
The next verification has a finite atomic acquired law, a singular physical law, expanding updates and changing mode masses. It provides a second realization of the general theorem, not a premise of that theorem.

Let $r_n\in[1/5,1/3]$, $\ell_0=1$, $\ell_n=\prod_{j=1}^n r_j$, and
\[
 x(\omega)=\sum_{n\ge1}(1-r_n)\ell_{n-1}\omega_n,
 \qquad \omega\in\{0,1\}^{\N}.
\]
The physical Cantor set $K$ is the image of this map. For a word $u$ of length $h$, let
\[
 c(u)=\sum_{n=1}^h(1-r_n)\ell_{n-1}u_n+\tfrac12\ell_h.
\]
Adjoin all these fair-tail means to $K$ to obtain $K^*$. It is compact: centers of unbounded depth approach $K$, and bounded depths give only finitely many centers. Finite centers represent posterior distributions, not additional physical Cantor points.

Prepare a mode $I_0\in\{0,1\}$ with any declared distribution, independently of a fair infinite tail. Report the mode and make $B$ paid calls revealing successive initial bits. At each of $T$ further opportunities, independently activate with any deterministic probability $\vartheta_t$. An inactive call copies both physical variables and costs one call. On activation, use the following raw stochastic operations, with $\gamma_0=1/4096$ and arbitrary deterministic $\alpha_t\in[0,1]$:
\begin{center}\small
\begin{tabular}{cccl}
\toprule
Source & Destination & Probability & Tail operation\\\midrule
$0$&$0$&$(1-\alpha_t)(1-\gamma_0)$&prepend two fair bits\\
$0$&$0$&$(1-\alpha_t)\gamma_0$&delete first bit\\
$0$&$1$&$\alpha_t$&delete first bit\\
$1$&$0$&$1/2$&prepend six fair bits\\
$1$&$1$&$1/2$&copy\\\bottomrule
\end{tabular}\end{center}
Report destination, operation and fresh prepended bits, but neither the deleted bit nor the old surviving tail. All kernels are nonnegative and normalized; operation probabilities are independent of bit values. Both mode changes transport a nonconstant old tail.

Let $w_T=\PP(I_T=1)$ under this actual protocol. The terminal Bernoulli probe has mean $X+s\delta$, where
\begin{equation}\label{eq:singulartask}
 X=\tfrac12+\tfrac14(I_T-w_T)+\tfrac18(x(\omega_T)-\tfrac12),
 \qquad s\in\{-1,1\},\quad0\le\delta\le1/16.
\end{equation}
These are valid means, and $\E X=1/2$. The mode and acquired cylinder posterior determine every legal continuation. Their score intervals are disjoint across modes and the current probe is strictly affine in the cylinder mean within a mode. They therefore realize the predictive quotient directly, including the declared time interface.

\begin{lemma}[Completed-domain geometry and raw gains]\label{lem:singularraw}
On $K^*$, deletion, fixing $c(\varnothing)$, is $30$-Lipschitz; prepending a fixed word of length $s$ is $6\,3^{-s}$-Lipschitz. On $S=\{0,1\}\times K^*$ put
\[
 d((i,x),(i,y))=\sqrt{v_i}|x-y|,
 \quad (v_0,v_1)=(1800,1),
 \quad d((0,x),(1,y))=2\sqrt{1800}.
\]
All centers of depth at most $h$ in both modes form an inward cover with
\[
 M_h=2^{h+2}-2,\qquad r_{M_h}=\sqrt{1800}\ell_h,\qquad w=1.
\]
The active true-report kernel satisfies the one-step pair inequality with $\kappa=17/20$, and its candidate envelope is identically one, independently of the mode law and schedule.
\end{lemma}
\begin{proof}
If two completed coordinates first differ, or one terminates, at level $k$, their separation lies between $\ell_{k-1}/6$ and $\ell_{k-1}$. Opposite branches are separated by $(1-2r_k)\ell_{k-1}\ge\ell_{k-1}/3$; a parent center and either branch are separated by $(1/2-r_k)\ell_{k-1}\ge\ell_{k-1}/6$. This also proves unique decoding of finite centers and infinite points. Deletion lowers the distinguishing level by one, giving gain at most $6/r_{k-1}\le30$. If the first digits differ or one point is the empty center, the gain is at most six. A common prepended word raises the distinguishing level by $s$, giving gain at most $6\ell_{k+s-1}/\ell_{k-1}\le6\,3^{-s}$.

The displayed metric satisfies the triangle inequality because within-mode diameters are at most $\sqrt{1800}$. Leave shorter centers fixed and truncate longer words to depth $h$. The cells are Borel, the error is at most $\sqrt{1800}\ell_h$, and the count includes both the mode and every possible retained word length.

For a candidate in the wrong mode, apply the reported operation to its coordinate and set its mode to the reported destination. This defines the common Borel update on every true/candidate input. For same-mode pairs the matrix of squared gains, with rows indexed by source and columns by destination, is
\[
 C_t=\begin{pmatrix}
 (1-\alpha_t)85/128&900\alpha_t\\
 2/59049&1/2
 \end{pmatrix}.
\]
Indeed $(1-1/4096)(4/9)+900/4096=85/128$. For $v=(1800,1)^{\mathsf T}$ both rows obey $(C_tv)_i\le(17/20)v_i$: the first is maximized at the larger of $1800(85/128)$ and $900$, and the second equals $3600/59049+1/2<17/20$. For cross-mode pairs the squared distance before the update is $4\cdot1800$ and afterward at most $1800$. The gain is at most $1/4$. These are inequalities under the physical report law for all pairs, not just a realized code trajectory. The constant envelope verifies the relation-robust conditions with $a=A=0$, $b=B=1$.
\end{proof}

\begin{theorem}[Actual-depth acquisition--memory law]\label{thm:singular}
Let $H_0=B$, and update the acquired depth by $h+2$, $h+6$, $\max(h-1,0)$ or $h$ according to the reported operation. Set $\omega_{T,i,h}=\PP(I_T=i,H_T=h)$. The actual nominal score law is
\begin{equation}\label{eq:atomic}
 \nu_{B,T}=\sum_{i,h}\omega_{T,i,h}2^{-h}
   \sum_{u\in\{0,1\}^h}
    \delta_{\,1/2+(i-w_T)/4+(c(u)-1/2)/8}.
\end{equation}
It is finite atomic. Put $A_{B,T}=\E(X-\E[X\mid\cH_{B,T}])^2$ and $b_\delta=\E[X(1-X)]-\delta^2$. For $M\ge6$ and $n=\lfloor\log_2 M\rfloor$, the sign-minimax risks satisfy
\begin{align}
 \mathcal R_\cp-b_\delta&=\delta^2+A_{B,T}+q_M(\nu_{B,T}),\label{eq:singularexact}\\
 \mathcal R_\cp-b_\delta&\asymp\mathcal R_\on-b_\delta
 \asymp\delta^2+\E\ell_{\min(H_T,n)}^2.\label{eq:singularlaw}
\end{align}
The constants are uniform over $B,T$, all ratios and schedules, and all initial mode laws. Raw acquisition is exactly $B+T+2$, including preparation, every inactive opportunity, and the terminal probe.
\end{theorem}
\begin{proof}
Inductively the full history knows a finite prefix, the unrevealed remainder is fair, and the prefix is uniform conditional on its depth and mode. Prepending discloses fresh independent bits, deletion reveals no unknown bit, and holds change nothing. This proves \eqref{eq:atomic}, with $H_T\le B+6T$. The unresolved physical coordinate variance in a depth-$h$ cylinder lies between $\ell_h^2/9$ and $\ell_h^2/4$: the first unresolved bit alone contributes at least $\ell_h^2/9$, and the whole cylinder has length $\ell_h$. Hence
\begin{equation}\label{eq:singularacquisition}
 \frac1{576}\E\ell_{H_T}^2\le A_{B,T}\le\frac1{256}\E\ell_{H_T}^2.
\end{equation}

Choose $h=\lfloor\log_2(M+2)\rfloor-2\ge1$. During the initial paid readings retain only the first $\min(B,h)$ bits. The seed is legal, has envelope one and error at most $\sqrt{1800}\ell_h$. Subsequently prepend and truncate, delete the first retained bit if present, or copy. Lemma~\ref{lem:singularraw}, Theorem~\ref{thm:main} with $m=1$, and Corollary~\ref{cor:holds} give nominal acquired-score distortion at most $C\ell_h^2\le C'\ell_n^2$, since $0\le n-h\le2$. The score is Lipschitz also across modes. The retained word is always a genuine prefix. Its origin pattern depends on the operations, not bit values; conditioning on that pattern and then averaging shows that the center readout is the conditional mean of $X$ given its retained word and mode. It is unbiased, so \eqref{eq:sign} adds exactly $\delta^2$.

For a converse, every operation preserves a fair physical tail conditional on the terminal mode. At least one terminal mode has mass at least $1/2$. Level-$k$ physical cylinders have length $\ell_k$ and gaps at least $\ell_k$. Their score gaps are at least $\ell_k/8$. A score ball of radius $\ell_k/32$ has diameter $\ell_k/16$, so meets at most one such cylinder of the selected mode. Its conditional mass is $2^{-k}$. Taking $k=n+2$ gives $2^{-k}\le1/(2M)$ and $\ell_k\ge\ell_n/25$. The union argument yields $q_M(\law(X))\ge c\ell_n^2$. This lower witness is the physical-oracle law; no small-ball regularity of the actual atomic law is asserted.

For every checkpoint forecast, averaging the two signs gives
\[
 \sup_s R_s-b_\delta\ge\delta^2+A_{B,T}+q_M(\nu_{B,T}).
\]
The finite atomic law has an optimal nominal code; replace its readouts by label-conditional means. This preserves its label count and makes its mean error zero, so
\[
 \sup_s R_s-b_\delta=\delta^2+A_{B,T}+q_M(\nu_{B,T})
\]
for such an optimal code. This proves \eqref{eq:singularexact}. Proposition~\ref{prop:acquired} and \eqref{eq:singularacquisition} compare both risks with $\delta^2+\E\ell_{H_T}^2+\ell_n^2$. Finally
\[
 \max\{\E\ell_{H_T}^2,\ell_n^2\}
 \le\E\ell_{\min(H_T,n)}^2
 \le\E\ell_{H_T}^2+\ell_n^2,
\]
which proves \eqref{eq:singularlaw}.
\end{proof}

This realization includes expanding deletion, arbitrary nonstationary mode masses and a nonhomogeneous sequence of singular scales. Its finite actual support changes with acquisition. The matching law is not obtained by assigning the physical Cantor dimension to the atomic acquired distribution. The known ratios may be supplied by a certified computable routine or a charged finite table to the required depth; arbitrary real constants do not by themselves furnish a finite program. The exact prefix update needs only the counted finite word, whereas numeric evaluation of the final center needs its own enclosure and program account.
